\chapter{System Implementation and Integrated Demonstration}

\section{Implementation Overview and Environment}
This chapter describes the implementation prototype of the multimodal target fusion and behavior analysis system. The purpose of this chapter is not to introduce a new algorithm, but to show how the system architecture defined in Chapter 2 and the methods developed in Chapters 3, 4, and 5 are organized into a runnable demonstration system.

The system environment is summarized in Table~\ref{tab:chapter6-env}. The implementation uses Python scripts for data processing, result adaptation, metric aggregation, and visualization payload generation. The demonstration interface is implemented as a static web page using HTML, CSS, and JavaScript. It can be opened locally in a browser or deployed as a static page for public demonstration. The experimental scripts are designed to use project-relative paths when transferring artifacts, so that the system is not bound to one fixed local directory.

\begin{table}[H]
  \centering
  \caption{Implementation environment of the demonstration system.}
  \label{tab:chapter6-env}
  \setlength{\tabcolsep}{4pt}
  \renewcommand{\arraystretch}{0.95}
  \small
  \begin{tabular}{p{0.24\linewidth}p{0.64\linewidth}}
    \toprule
    Component & Implementation description \\
    \midrule
    Offline processing & Python-based scripts for reading datasets, adapting experimental outputs, computing metrics, and generating visualization records. \\
    Data artifacts & Sequence metadata, trajectory records, anomaly labels or protocols, method scores, registration diagnostics, and frame-level visualization payloads. \\
    Demonstration interface & Static browser-based page implemented with HTML, CSS, and JavaScript. \\
    Deployment mode & Local browser preview and static-page deployment, such as GitHub Pages. \\
    Path organization & Relative project paths are used for generated artifacts to improve portability across local and remote environments. \\
    \bottomrule
  \end{tabular}
\end{table}

The implemented system follows the six functional layers introduced in Chapter 2: multimodal data input, preprocessing, optical reference retrieval, entity-level fusion, behavior analysis, and result presentation. In the prototype, algorithmic experiments and result generation are performed offline. The system then organizes the generated trajectories, anomaly scores, registration diagnostics, heatmaps, and explanation records into a unified visualization package. This design allows the experimental algorithms to be updated independently while keeping the presentation interface stable.

This environment summary defines the engineering boundary of the prototype. The following module description then explains how the implementation corresponds to the architectural layers introduced earlier.

\section{Open-source Repository and Reproducibility}
The implementation code, thesis source files, experimental scripts, and visualization materials of this project are organized in a public GitHub repository: \url{https://github.com/wcx12/FusionTrack}. The repository is used as the reproducibility entry for the implemented system. It provides the LaTeX source of the thesis, anomaly detection benchmark scripts, VPR and registration modules, static visualization materials, and workflow files for document building.

The open-source package is organized around the same modular boundary as the system implementation. The benchmark code defines data protocols, result adapters, metric computation, and key-audit rules. The visualization materials provide the static demonstration page and generated result files used for thesis presentation. The repository does not include large raw datasets, server credentials, private configuration files, or absolute local paths. This separation keeps the public materials inspectable while avoiding dependence on one local machine or remote server account.

\section{Implementation of Main Modules}
To keep the implementation consistent with the system model in Chapter 2, the prototype maps each architectural layer to a concrete system module. The data input and preprocessing modules form the entry of the system. They read the available dataset sequences and experimental results, then normalize them into records indexed by sequence, frame, target, method, and score. This organization follows the target-centered representation defined in Chapter 2 and makes later result comparison more direct.

The optical reference retrieval module corresponds to Chapter 3. In the system implementation, TF-VPR is treated as a reference retrieval component for coordinate-missing optical images. Its output is used as a scene-level constraint or as an experimental result interface, rather than as a direct target identity decision.

The entity-level fusion and registration module corresponds to Chapter 4. The implementation preserves the registration and fusion outputs in a form suitable for inspection, including correspondence-related records, fused trajectory information, and diagnostic indicators. These outputs support the later behavior analysis module.

The behavior analysis module corresponds to Chapter 5. It organizes individual anomaly results, group anomaly results, fused anomaly evidence, event explanations, and situation heatmaps. The module does not only output a scalar score; it also preserves the visual and textual evidence required for interpreting abnormal behavior.

The result presentation module finally converts these records into an interactive demonstration interface. It presents metrics, dynamic visualization, method comparison, explanation records, language switching, and report export, so that the outputs of the previous modules can be inspected in a unified page.

\section{System Workflow}
The runtime workflow of the implemented system is organized around experimental result loading and visualization generation. First, the system reads dataset metadata, frame resources, trajectory records, and method outputs. Second, the result adaptation module converts these files into a unified format. Third, the metric module summarizes sequence-level and method-level indicators. Fourth, the visualization module generates the dynamic payload used by the web interface, including original frames, trajectory overlays, heatmaps, and combined views. Finally, the static page loads these records and provides interactive inspection.

In compact form, the workflow consists of data and result loading, unified result adaptation, metric aggregation, visualization payload generation, and web-based demonstration with export.

This workflow separates algorithm execution from result presentation. Therefore, when a stronger detector, fusion method, or registration algorithm is added later, it only needs to provide outputs in the agreed result format. The visualization layer can then reuse the same interface to display the updated results.

\section{Integrated Demonstration Interface}
The integrated demonstration interface is designed to support thesis presentation and result inspection. The top part of the interface provides method selection, sequence selection, language switching, and core indicators. These controls allow the user to compare different methods and inspect the results of different sequences without changing the underlying data files.

The central visualization area provides four synchronized views: the original video, the anomaly heatmap, the trajectory view, and the combined heatmap-trajectory view. The original video keeps the scene context visible. The heatmap view shows where abnormal evidence is concentrated. The trajectory view highlights target movement patterns. The combined view is used as the default display because it presents spatial risk and target motion at the same time.

The system also separates different result types. Individual-level results focus on abnormal behavior of a single target, such as route deviation or motion irregularity. Group-level results focus on relational changes among multiple targets, such as leaving, merging, splitting, or neighborhood restructuring. Registration diagnostics are presented separately to avoid mixing cross-modal alignment quality with video anomaly detection results.

In addition to dynamic visualization, the interface provides explanation records, protocol information, data audit information, and result export. These functions make the system more suitable for thesis demonstration: the main page can remain concise during presentation, while detailed evidence can still be expanded when needed.

\section{Chapter Summary}
This chapter presented the implementation prototype and integrated demonstration of the proposed system. Different from Chapter 2, which defines the system architecture and data model, this chapter focuses on how the architecture is realized in an offline processing and web-based presentation workflow.

The implemented prototype connects the TF-VPR reference retrieval interface, entity-level fusion and registration outputs, individual and group anomaly analysis results, and situation heatmap visualization. Through unified result adaptation and static web presentation, the system provides a runnable demonstration for the methods developed in the previous chapters. The prototype also keeps the result interface modular, allowing later algorithmic improvements to be integrated without redesigning the whole visualization system.
